\section{Discussion}
\label{sec:discussion}

\paragraph{Low-dimensional behavior.}
At $d = 5$, \mmr shows worse uniformity (CV $\approx 2.2$) than Random Search (CV $\approx 0.21$).
With $K = 1000$ solutions in a 5D unit hypercube, even random points achieve good coverage; the fitness term ($\lambda = 0.5$) then creates clusters around fitness optima.
This is not a failure of the MMR formulation but a regime where diversity pressure is unnecessary---analogous to how regularization hurts when the model is already underfitting.
Practitioners working in low-dimensional behavior spaces ($d \leq 5$) may prefer standard MAP-Elites, which provides excellent coverage with manageable grid sizes ($3^5 = 243$ cells).

\paragraph{Saturating distances.}
The exponential saturating distance (Equation~\ref{eq:saturating}) addresses the boundary effect of Euclidean distance, where unbounded growth rewards extreme separation over interior coverage.
By saturating at ``different enough,'' these functions promote uniform interior filling.
The scale parameter $\sigma$ can be set automatically from the data distribution, avoiding manual tuning.

\paragraph{When to use \mmr vs.\ MAP-Elites.}
We recommend \mmr when: (1) the behavior space exceeds 5 dimensions; (2) the behavior space lacks natural axes for discretization; (3) non-Euclidean distances are appropriate; or (4) a fixed-size archive is required.
MAP-Elites remains preferable for low-dimensional behavior spaces where its grid provides interpretable coverage guarantees.

\paragraph{Connections to submodular optimization.}
The diversity component of \mmr's objective is related to the facility location problem and determinantal point processes~\cite{kulesza2012determinantal}.
While the full fitness-diversity objective is not submodular, the greedy approach inherits strong empirical performance from this connection.
Future work could explore tighter approximation bounds for the mixed objective.

\paragraph{Limitations.}
\mmr requires a meaningful distance function over the behavior space.
When behaviors lack natural similarity structure, the algorithm reduces to fitness-only selection.
The $\lambda$ parameter introduces a hyperparameter not present in MAP-Elites, though we find performance robust across $\lambda \in [0.3, 0.7]$.
The current implementation uses Gaussian mutation; integration with advanced emitters (CMA-ES, policy gradients) is straightforward but not yet explored.
